% !PRINTED 24
\weberpage{24}
%<<CONT
that the index series $j_{r+1}, j_{r+2} \ldots j_\mu$ and $j_{r+1}',
j_{r+2}' \ldots j_\mu'$ agree, apart from the order.
Consequently the preceding parts of the index series
\[
j_1, j_2 \ldots j_r \quad \text{and} \quad j_1', j_2' \ldots j_r'
\]
must also agree.

\begin{center}
\S.\,7.

\gesperrt{Further Theorems on the Composition Series.}
\end{center}

Of the composition series many further theorems hold, of
which we here derive the most important.

\gesperrt{II. If $P_\nu$ is any normal divisor of $P$, then there
is a composition series of $P$ in which $P_\nu$ occurs.}

If $P_\nu$ is a greatest normal divisor of $P$, we can set $P_\nu =
P_1$ and continue the composition series from there on arbitrarily.
But if $P_\nu$ is not a greatest normal divisor of $P$, then we
seek a greatest normal divisor above $P_\nu$, which we denote by
$P_1$. Then $P_\nu$ is a normal divisor of $P_1$. If it is a
greatest one, we set $P_\nu = P_2$ and continue the composition
series from there arbitrarily. But if $P_\nu$ is still not a greatest
normal divisor of $P_1$, we seek a greatest normal divisor of
$P_1$ above $P_\nu$ and proceed in this way until we finally
arrive at $P_\nu$, which must necessarily occur after a finite number
of steps. If we then append a composition series of $P_\nu$, we
have a composition series of $P$ in which $P_\nu$ occurs, as
theorem II. requires.

In a composition series of a group $P$, by virtue of the
definition, every term is a normal divisor of every preceding one.
The first term $P_1$ and the last term 1 are at the same time
normal divisors of $P$ itself. For the intermediate $P_2, P_3
\ldots$ this will in general not be the case. In the composition
series there are thus certain distinguished terms which are at the
same time normal divisors of $P$ itself, and which play a special
role\footnote{They form the so-called principal series of $P$
(\gesperrt{Netto}).}. On this we derive in what follows one more
important theorem.

Suppose $Q$ is a term of a composition series of $P$ which is
at the same time a normal divisor of $P$, while the term following upon $Q$,
%CONT>>
